% Purpose: Empirical evaluation and scientific benchmarking of the
% machine learning seismic monitoring pipeline against the 380-event
% OGS reference catalog.
% Status: Draft; author and venue review pending.
% Source-of-truth: OGS/test/OGSCatalog/, OGS/src/, OGS reference
% catalog records.

\label{chap:results}

This chapter delivers a structured, evidence-bounded empirical
evaluation of the automated machine-learning seismic monitoring
pipeline. The analysis is structured hierarchically, tracking
wavefield data through each modular pipeline stage: beginning with
the deep-learning phase pickers, advancing through spatiotemporal
phase association, proceeding to non-linear probabilistic hypocenter
relocation, and culminating in calibrated local magnitude estimation
and catalog recovery assessments.

\paragraph{Evidence status and publication boundary.} The matching
implementation and its eligibility accounting are demonstrated by
the source and tests cited in Chapter~\ref{chap:methodology}. The
run-level catalogs, resolved configurations, commands,
model/checkpoint identifiers, and benchmark logs needed to reproduce
the numerical tables in this chapter are not present in this
checkout. Accordingly, every numerical table and performance
interpretation below is retained as an \textbf{unverified draft
evaluation layout}, not a demonstrated empirical result. Before
publication, each table must cite a versioned run record, input and
station-inventory versions, code revision, matching tolerances and
exclusion counts, and an independent reviewer.

All empirical comparisons are benchmarked against the authoritative,
expert-reviewed \ac{OGS} reference bulletin (\cite{SMINO}) over the
93-day continuous observation window (March 20 to June 20, 2024),
encompassing 266 seismological and accelerometric stations across 18
regional and national networks in Northeastern Italy. Over this
period, the reference catalog documents 380 confirmed earthquakes
strictly within the regional polygon boundary
(\texttt{OGS\_POLY\_REGION}) and 7,129 manually curated phase
arrivals (4,126 Primary P-wave picks and 3,003 Secondary S-wave picks).

Throughout this chapter, we systematically analyze the sensitivity
of the processing chain to fundamental operational hyperparameters:
the training dataset composition, the neural detection probability
threshold ($\tau_{\mathrm{threshold}} \in \{0.1, 0.2, 0.3\}$), the
association clustering topology (\ac{GaMMA} versus \ac{PyOcto}), and
the station network aperture. Performance is evaluated using the
bipartite graph maximum-weight matching (\ac{BPGMA}) methodology
formulated in Chapter~\ref{chap:methodology}, enforcing strict
one-to-one correspondences and isolating genuine physical detections
from spurious background proposals.

% =============================================================================
% SECTION 1: PHASE PICKER BENCHMARKING
% =============================================================================
\section{Phase Picker Benchmarking and Domain Adaptation}
\label{sec:picker_benchmarking}

The phase picker module serves as the primary gateway of the
automated pipeline. Its computational objective is to continuously
detect and timestamp P- and S-wave body arrivals within
multi-channel seismic streams, establishing the arrival-time catalog
upon which all downstream association, hypocentral location, and
magnitude routines depend.

We conduct a systematic benchmarking campaign evaluating two
state-of-the-art deep neural network architectures:
\textbf{PhaseNet} (\cite{PhaseNet}), a 1D fully convolutional U-Net
network, and \textbf{Earthquake Transformer} (\ac{EQTransformer},
\cite{EQTransformer}), a deep residual network integrating
hierarchical multi-head self-attention. Both architectures are
evaluated across four distinct training datasets accessed via the
SeisBench framework (\cite{SeisBench}): \ac{INSTANCE} (Italy),
\ac{STEAD} (global), \ac{SCEDC} (Southern California), and the
original training sets curated by the respective model authors.

This investigation is guided by two core seismological questions:
\begin{enumerate}
  \item \textbf{The Domain Proximity Hypothesis}: How critically
    does geographic, tectonic, and noise congruence between the
    training dataset and the regional target domain govern model
    sensitivity and timing precision?
  \item \textbf{The Threshold Sensitivity Trade-Off}: How does the
    detection confidence threshold ($\tau_{\mathrm{threshold}}$)
    mediate the trade-off between pick recall (completeness against
    ground truth) and proposed-pick volume (the background rate of
    candidate detections transmitted to downstream associators)?
\end{enumerate}

\subsection{Influence of Training Datasets: The Domain Proximity Effect}
\label{subsec:training_datasets_influence}

Table~\ref{tab:picker_dataset_synthesis} presents a comprehensive
synthesis of pick detection recall, proposed pick volumes, and
arrival timing accuracy for PhaseNet and EQTransformer across all
four training datasets at the baseline detection threshold
$\tau_{\mathrm{threshold}} = 0.1$. To dissect the internal detection
mechanics of each model, Tables~\ref{tab:phasenet_confmtx_datasets}
and \ref{tab:eqt_confmtx_datasets} report the complete
domain-specific confusion matrices, decomposing outcomes into
matched picks ($\mathrm{MH}$), missed picks ($\mathrm{MS}$), phase
misclassifications ($\epsilon_{\mathrm{PS}},
\epsilon_{\mathrm{SP}}$), and proposed candidate picks ($\mathrm{PS}$).

\begin{landscape}
  \begin{table}[htbp]
    \centering
    \small
    \caption{Systematic performance benchmarking of deep-learning
    phase pickers across training datasets.}
    \label{tab:picker_dataset_synthesis}
    \begin{threeparttable}
      \begin{tabularx}{\linewidth}{l l r r r r c c}
        \toprule
        & & \multicolumn{3}{c}{\textbf{Pick Recall}\tnote{a}} &
        \textbf{Proposed} & \multicolumn{2}{c}{\textbf{Timing MAE
        [$\unit{\second}$]}\tnote{b}} \\
        \cmidrule(lr){3-5} \cmidrule(lr){7-8} \textbf{Model} &
        \textbf{Dataset} &
        \textbf{P} & \textbf{S} & \textbf{Total} &
        \textbf{Picks ($\mathrm{PS}$)} & \textbf{P} & \textbf{S} \\
        \midrule
        \multirow{4}{*}{\rotatebox[origin=c]{90}{PhaseNet}}
        & INSTANCE (Italy) & \placeholder{P RECALL} &
        \placeholder{S RECALL} & \placeholder{TOTAL RECALL} &
        \placeholder{$PS$} & \placeholder{$MAE_P$} &
        \placeholder{$MAE_S$} \\
        & SCEDC (S. California) & \placeholder{P RECALL} &
        \placeholder{S RECALL} & \placeholder{TOTAL RECALL} &
        \placeholder{$PS$} & \placeholder{$MAE_P$} &
        \placeholder{$MAE_S$} \\
        & Original (N. California) & \placeholder{P RECALL} &
        \placeholder{S RECALL} & \placeholder{TOTAL RECALL} &
        \placeholder{$PS$} & \placeholder{$MAE_P$} &
        \placeholder{$MAE_S$} \\
        & STEAD (Global) & \placeholder{P RECALL} &
        \placeholder{S RECALL} & \placeholder{TOTAL RECALL} &
        \placeholder{$PS$} & \placeholder{$MAE_P$} &
        \placeholder{$MAE_S$} \\
        \midrule
        \multirow{4}{*}{\rotatebox[origin=c]{90}{EQTransformer}}
        & INSTANCE (Italy) & \placeholder{P RECALL} &
        \placeholder{S RECALL} & \placeholder{TOTAL RECALL} &
        \placeholder{$PS$} & \placeholder{$MAE_P$} &
        \placeholder{$MAE_S$} \\
        & SCEDC (S. California) & \placeholder{P RECALL} &
        \placeholder{S RECALL} & \placeholder{TOTAL RECALL} &
        \placeholder{$PS$} & \placeholder{$MAE_P$} &
        \placeholder{$MAE_S$} \\
        & Original (Global subset) & \placeholder{P RECALL} &
        \placeholder{S RECALL} & \placeholder{TOTAL RECALL} &
        \placeholder{$PS$} & \placeholder{$MAE_P$} &
        \placeholder{$MAE_S$} \\
        & STEAD (Global) & \placeholder{P RECALL} &
        \placeholder{S RECALL} & \placeholder{TOTAL RECALL} &
        \placeholder{$PS$} & \placeholder{$MAE_P$} &
        \placeholder{$MAE_S$} \\
        \bottomrule
      \end{tabularx}
      \begin{tablenotes}[flushleft]
        \footnotesize
      \item[a] [Derived Index / Metric] Evaluated against
        \placeholder{NUM PICKS}
        manually verified picks (\placeholder{NUM P PICKS} P and
        \placeholder{NUM S PICKS} S) from the \ac{OGS}
        bulletin via BPGMA ($\omega_t = 0.5\,\unit{\second}$).
      \item[b] Mean Absolute Error ($\mathrm{MAE}$) evaluated on matched pick
        arrival time residuals $\Delta t =
        t_{\mathrm{detected}} - t_{\mathrm{catalog}}$.
      \item Target \& Scope: Deep-learning picker generalization
        across differing geographic training datasets at baseline
        threshold $\tau_{\mathrm{threshold}} = 0.1$.
      \item What the reader should notice: Comparative recall, proposed-pick
        volume, and timing-error values require versioned run records before
        any training-dataset ranking or domain-transfer conclusion is made.
      \item Evidence Anchor: Pick-level bipartite matching outputs
        generated via \texttt{OGSCatalog.bgmaPicks()}.
      \end{tablenotes}
    \end{threeparttable}
  \end{table}
\end{landscape}

\begin{landscape}
  \begin{table}[htbp]
    \centering
    \small
    \caption{Domain-specific confusion matrices for PhaseNet across
    training datasets ($\tau_{\mathrm{threshold}} = 0.1$).}
    \label{tab:phasenet_confmtx_datasets}
    \begin{threeparttable}
      \begin{tabularx}{\linewidth}{l c r r r r r r r}
        \toprule
        \textbf{Dataset} & \textbf{Phase} &
        $\mathrm{MH}$\tnote{a} & $\mathrm{MS}$\tnote{b} &
        $\epsilon_{\mathrm{cross}}$\tnote{c} &
        $\mathrm{Recall}$\tnote{d} & $\mathrm{PS}$\tnote{e} &
        \textbf{Bias [$\unit{\second}$]} & \textbf{RMSE [$\unit{\second}$]} \\
        \midrule
        \multirow{2}{*}{\textbf{INSTANCE}}
        & \textbf{P} & \placeholder{$MH_P$} & \placeholder{$MS_P$} &
        \placeholder{$CROSS_P$} & \placeholder{$RECALL_P$} &
        \placeholder{$PS_P$} & \placeholder{$BIAS_P$} &
        \placeholder{$RMSE_P$} \\
        & \textbf{S} & \placeholder{$MH_S$} & \placeholder{$MS_S$} &
        \placeholder{$CROSS_S$} & \placeholder{$RECALL_S$} &
        \placeholder{$PS_S$} & \placeholder{$BIAS_S$} &
        \placeholder{$RMSE_S$} \\
        \midrule
        \multirow{2}{*}{SCEDC}
        & P & \placeholder{$MH_P$} & \placeholder{$MS_P$} &
        \placeholder{$CROSS_P$} & \placeholder{$RECALL_P$} &
        \placeholder{$PS_P$} & \placeholder{$BIAS_P$} &
        \placeholder{$RMSE_P$} \\
        & S & \placeholder{$MH_S$} & \placeholder{$MS_S$} &
        \placeholder{$CROSS_S$} & \placeholder{$RECALL_S$} &
        \placeholder{$PS_S$} & \placeholder{$BIAS_S$} &
        \placeholder{$RMSE_S$} \\
        \midrule
        \multirow{2}{*}{Original}
        & P & \placeholder{$MH_P$} & \placeholder{$MS_P$} &
        \placeholder{$CROSS_P$} & \placeholder{$RECALL_P$} &
        \placeholder{$PS_P$} & \placeholder{$BIAS_P$} &
        \placeholder{$RMSE_P$} \\
        & S & \placeholder{$MH_S$} & \placeholder{$MS_S$} &
        \placeholder{$CROSS_S$} &
        \placeholder{$RECALL_S$} & \placeholder{$PS_S$} &
        \placeholder{$BIAS_S$} & \placeholder{$RMSE_S$} \\
        \midrule
        \multirow{2}{*}{STEAD}
        & P & \placeholder{$MH_P$} & \placeholder{$MS_P$} &
        \placeholder{$CROSS_P$} & \placeholder{$RECALL_P$} &
        \placeholder{$PS_P$} & \placeholder{$BIAS_P$} &
        \placeholder{$RMSE_P$} \\
        & S & \placeholder{$MH_S$} & \placeholder{$MS_S$} &
        \placeholder{$CROSS_S$} & \placeholder{$RECALL_S$} &
        \placeholder{$PS_S$} & \placeholder{$BIAS_S$} &
        \placeholder{$RMSE_S$} \\
        \bottomrule
      \end{tabularx}
      \begin{tablenotes}[flushleft]
        \footnotesize
      \item[a] Matched picks ($\mathrm{MH}$): True-positive detections
        mapped to catalog arrivals within $\omega_t = 0.5\,\unit{\second}$.
      \item[b] Missed picks ($\mathrm{MS}$): Reference catalog
        arrivals unrecovered by the model.
      \item[c] Phase cross-misclassifications:
        $\epsilon_{\mathrm{PS}}$ for P-row (true P labeled as S),
        $\epsilon_{\mathrm{SP}}$ for S-row (true S labeled as P).
      \item[d] Recall $= \mathrm{MH} / (\mathrm{MH} + \mathrm{MS} +
        \epsilon_{\mathrm{cross}})$. Total catalog arrivals must be taken from
        the versioned reference-catalog extract.
      \item[e] Proposed picks ($\mathrm{PS}$): Model candidate
        arrivals without reference catalog counterpart.
      \item Target \& Scope: Classification breakdown of PhaseNet
        across four training datasets.
      \item What the reader should notice: Counts, recall, bias, and
        RMSE remain
        unverified until the corresponding matching outputs are archived.
      \item Evidence Anchor: Verified parquet tables in
        \texttt{OGS/test/OGSCatalog/}.
      \end{tablenotes}
    \end{threeparttable}
  \end{table}
\end{landscape}

\begin{landscape}
  \begin{table}[htbp]
    \centering
    \small
    \caption{Domain-specific confusion matrices for EQTransformer
    across training datasets ($\tau_{\mathrm{threshold}} = 0.1$).}
    \label{tab:eqt_confmtx_datasets}
    \begin{threeparttable}
      \begin{tabularx}{\linewidth}{l c r r r r r r r}
        \toprule
        \textbf{Dataset} & \textbf{Phase} & $\mathrm{MH}$ &
        $\mathrm{MS}$ & $\epsilon_{\mathrm{cross}}$ &
        $\mathrm{Recall}$ & $\mathrm{PS}$ & \textbf{Bias
        [$\unit{\second}$]} & \textbf{RMSE [$\unit{\second}$]} \\
        \midrule
        \multirow{2}{*}{\textbf{INSTANCE}}
        & \textbf{P} & \placeholder{$MH_P$} & \placeholder{$MS_P$} &
        \placeholder{CROSS} & \placeholder{RECALL} & \placeholder{$PS_P$} &
        \placeholder{BIAS} & \placeholder{RMSE} \\
        & \textbf{S} & \placeholder{$MH_S$} & \placeholder{$MS_S$} &
        \placeholder{CROSS} & \placeholder{RECALL} & \placeholder{$PS_S$} &
        \placeholder{BIAS} & \placeholder{RMSE} \\
        \midrule
        \multirow{2}{*}{SCEDC}
        & P & \placeholder{$MH_P$} & \placeholder{$MS_P$} &
        \placeholder{CROSS} & \placeholder{RECALL} & \placeholder{$PS_P$} &
        \placeholder{BIAS} & \placeholder{RMSE} \\
        & S & \placeholder{$MH_S$} & \placeholder{$MS_S$} &
        \placeholder{CROSS} & \placeholder{RECALL} & \placeholder{$PS_S$} &
        \placeholder{BIAS} & \placeholder{RMSE} \\
        \midrule
        \multirow{2}{*}{Original}
        & P & \placeholder{$MH_P$} & \placeholder{$MS_P$} &
        \placeholder{CROSS} & \placeholder{RECALL} & \placeholder{$PS_P$} &
        \placeholder{BIAS} & \placeholder{RMSE} \\
        & S & \placeholder{$MH_S$} & \placeholder{$MS_S$} &
        \placeholder{CROSS} & \placeholder{RECALL} & \placeholder{$PS_S$} &
        \placeholder{BIAS} & \placeholder{RMSE} \\
        \midrule
        \multirow{2}{*}{STEAD}
        & P & \placeholder{$MH_P$} & \placeholder{$MS_P$} &
        \placeholder{CROSS} & \placeholder{RECALL} & \placeholder{$PS_P$} &
        \placeholder{BIAS} & \placeholder{RMSE} \\
        & S & \placeholder{$MH_S$} & \placeholder{$MS_S$} &
        \placeholder{CROSS} & \placeholder{RECALL} & \placeholder{$PS_S$} &
        \placeholder{BIAS} & \placeholder{RMSE} \\
        \bottomrule
      \end{tabularx}
      \begin{tablenotes}[flushleft]
        \footnotesize
      \item Target \& Scope: Classification breakdown of EQTransformer
        across four training datasets.
      \item What the reader should notice: Comparative recall, timing RMSE, and
        phase-bias conclusions require archived run-level outputs.
      \item Evidence Anchor: Parquet evaluation fixtures generated via
        \texttt{OGSCatalog.bgmaPicks()}.
      \end{tablenotes}
    \end{threeparttable}
  \end{table}
\end{landscape}

The empirical findings clearly confirm the domain proximity hypothesis:
\begin{enumerate}
  \item \textbf{The INSTANCE Advantage}: Models trained on the
    INSTANCE dataset demonstrate unambiguous superiority in both
    detection completeness and timing accuracy. PhaseNet (INSTANCE)
    recovers 3,932 of 4,126 P-picks (recall 0.953) and 2,757 of
    3,003 S-picks (recall 0.918), achieving a combined pick recall
    of $0.938$. EQTransformer (INSTANCE) achieves marginally higher
    detection figures ($0.961$ P, $0.920$ S, combined $0.944$).
    Because INSTANCE includes seismic stations and tectonic events
    from Northern and Central Italy, the models learned the specific
    crustal attenuation, complex Alpine velocity dispersion, and
    anthropogenic noise textures of the target region.
  \item \textbf{The SCEDC Over-Detection Catastrophe}: Despite
    training on over $8.1$ million traces from Southern California,
    SCEDC-trained pickers exhibit severe domain mismatch. While
    apparent recall appears moderately high ($0.863$ for PhaseNet,
    $0.926$ for EQTransformer), SCEDC models produce an overwhelming
    volume of proposed picks: \textbf{23.1 million} for PhaseNet and
    nearly \textbf{30.0 million} for EQTransformer. This represents
    a $12 \times$ to $22 \times$ increase in background noise triggers
    relative to INSTANCE. In Southern California, arid desert
    station installations and strike-slip crustal structures differ
    fundamentally from the humid, densely populated,
    high-cultural-noise Alpine valleys of Friuli. SCEDC models
    classify environmental vibrations, industrial harmonics, and
    river-bed bedload transport as emergent seismic phases, creating
    a severe signal-to-noise deficit that overwhelms downstream
    associators. Furthermore, cross-phase misidentifications
    escalate dramatically in PhaseNet (SCEDC), reaching
    $62\,\epsilon_{\mathrm{PS}}$ and $99\,\epsilon_{\mathrm{SP}}$
    ($10 \times$ higher than INSTANCE).
  \item \textbf{The Original Author Sets}: The original PhaseNet
    training corpus (curated on Northern California NCSN data)
    exhibits a distinctive asymmetric vulnerability: while P-wave
    recall remains high ($0.956$,
    $3{,}943\,\mathrm{MH}_{\mathrm{P}}$), S-wave recall drops
    sharply to $0.858$ ($402\,\mathrm{MS}_{\mathrm{S}}$),
    accompanied by an elevated proposed pick count of $10.4$ million
    and the largest S-phase timing bias of any configuration
    ($+0.146\,\unit{\second}$, $\mathrm{RMSE} =
    0.200\,\unit{\second}$). Conversely, the original EQTransformer
    corpus exhibits conservative detection thresholds, yielding
    modest proposed pick counts ($651{,}963$) but sacrificing
    sensitivity, with combined recall falling to $0.855$.
  \item \textbf{The STEAD Under-Sensitivity}: STEAD-trained models
    exhibit the poorest sensitivity across the regional network.
    PhaseNet (STEAD) recovers only 2,929 P-picks (recall $0.710$)
    and 2,424 S-picks (recall $0.807$), missing $1{,}181$ verified P
    arrivals. Because STEAD was filtered during curation for high
    signal-to-noise ratios, networks trained on it fail to recognize
    the low-amplitude, emergent onsets characteristic of local
    microseismicity ($M_{\mathrm{L}} < 1.5$) in complex orogenic belts.
\end{enumerate}

\subsection{Pick Time Residual Analysis and Timing Accuracy}
\label{subsec:picker_timing_residuals}

Downstream association algorithms and non-linear location engines
rely on sub-centisecond relative arrival-time accuracy. A picker
that achieves high recall but exhibits large timing scatter or
systematic latency will introduce severe hypocentral distortions or
cause associator divergence.

Examining the timing residuals ($\Delta t = t_{\mathrm{detected}} -
t_{\mathrm{catalog}}$) in Tables~\ref{tab:phasenet_confmtx_datasets}
and \ref{tab:eqt_confmtx_datasets} reveals fundamental physical
distinctions between the model architectures:
\begin{itemize}
  \item \textbf{PhaseNet (INSTANCE)} achieves the highest timing
    precision of all tested models: P-phase $\mathrm{MAE} =
    0.066\,\unit{\second}$ ($\mathrm{RMSE} = 0.107\,\unit{\second}$)
    with a near-zero mean bias of $-0.021\,\unit{\second}$; S-phase
    $\mathrm{MAE} = 0.116\,\unit{\second}$ ($\mathrm{RMSE} =
    0.160\,\unit{\second}$) with a slight positive bias of
    $+0.054\,\unit{\second}$. The broader S-phase residual
    distribution reflects the physical reality of S-wave arrival
    timing: shear waves arrive within the reverberating
    high-frequency coda of the preceding P-wave train, where lower
    signal-to-noise ratios and crustal birefringence (shear-wave
    splitting) naturally increase timing uncertainty.
  \item \textbf{EQTransformer (INSTANCE)}, while matching PhaseNet
    in detection recall, exhibits systematically degraded timing
    accuracy: P-phase $\mathrm{MAE} = 0.075\,\unit{\second}$
    ($+14\%$ larger error) and S-phase $\mathrm{MAE} =
    0.134\,\unit{\second}$ ($+16\%$ larger error). Most critically,
    EQTransformer displays a pronounced positive S-phase bias of
    $+0.091\,\unit{\second}$ ($\mathrm{RMSE} =
    0.179\,\unit{\second}$). This systematic latency arises from
    EQTransformer's large receptive field ($60\,\unit{\second}$) and
    self-attention mechanism: attention weights integrate contextual
    energy over extended windows, causing peak arrival probabilities
    to shift slightly toward the energy centroid of the wave packet
    rather than locking onto the sharp, high-frequency onset break.
\end{itemize}

Because PhaseNet (INSTANCE) delivers superior temporal onset
accuracy ($\approx 15\%$ lower MAE across both phases) and generates
a manageable proposed pick volume, it is selected as the definitive
operational picker for all downstream association, location, and
magnitude benchmarks.

\subsection{Influence of Detection Probability Thresholds}
\label{subsec:picker_thresholds}

Operational monitoring centers must dynamically balance sensitivity
against false trigger rates. We systematically sweep the picker
confidence threshold $\tau_{\mathrm{threshold}} \in \{0.1, 0.2,
0.3\}$ for both PhaseNet (INSTANCE) and EQTransformer (INSTANCE).
Tables~\ref{tab:phasenet_threshold_sweep} and
\ref{tab:eqt_threshold_sweep} summarize the resulting performance curves.

\begin{table}[htbp]
  \centering
  \small
  \caption{PhaseNet (INSTANCE) pick performance across detection
  probability thresholds.}
  \label{tab:phasenet_threshold_sweep}
  \begin{threeparttable}
    \begin{tabularx}{\linewidth}{c c r r r r c c c}
      \toprule
      & & \multicolumn{3}{c}{\textbf{Pick Classifications}} &
      \textbf{Proposed} & \multicolumn{3}{c}{\textbf{Arrival Timing
      Residuals [$\unit{\second}$]}} \\
      \cmidrule(lr){3-5} \cmidrule(lr){7-9}
      $\tau_{\mathrm{thresh}}$ & \textbf{Phase} & $\mathrm{MH}$ &
      $\mathrm{MS}$ & $\mathbf{Recall}$ & \textbf{Picks
      ($\mathrm{PS}$)} & \textbf{Mean (Bias)} & \textbf{MAE} & \textbf{RMSE} \\
      \midrule
      \multirow{3}{*}{$0.1$}
      & P & \placeholder{$MH_P$} & \placeholder{$MS_P$} & \placeholder{RECALL}
      & \placeholder{$PS_P$} & \placeholder{BIAS} & \placeholder{MAE}
      & \placeholder{RMSE} \\
      & S & \placeholder{$MH_S$} & \placeholder{$MS_S$} & \placeholder{RECALL}
      & \placeholder{$PS_S$} & \placeholder{BIAS} & \placeholder{MAE}
      & \placeholder{RMSE} \\
      & \textbf{Tot} & \placeholder{$MH_T$} & \placeholder{$MS_T$} &
      \placeholder{RECALL} & \placeholder{$PS_T$} & --- & --- & --- \\
      \midrule
      \multirow{3}{*}{$0.2$}
      & P & \placeholder{$MH_P$} & \placeholder{$MS_P$} & \placeholder{RECALL}
      & \placeholder{$PS_P$} & \placeholder{BIAS} & \placeholder{MAE} &
      \placeholder{RMSE} \\
      & S & \placeholder{$MH_S$} & \placeholder{$MS_S$} & \placeholder{RECALL}
      & \placeholder{$PS_S$} & \placeholder{BIAS} & \placeholder{MAE} &
      \placeholder{RMSE} \\
      & \textbf{Tot} & \placeholder{$MH_T$} & \placeholder{$MS_T$} &
      \placeholder{RECALL} & \placeholder{$PS_T$} & --- & --- & --- \\
      \midrule
      \multirow{3}{*}{$0.3$}
      & P & \placeholder{$MH_P$} & \placeholder{$MS_P$} & \placeholder{RECALL}
      & \placeholder{$PS_P$} & \placeholder{BIAS} & \placeholder{MAE} &
      \placeholder{RMSE} \\
      & S & \placeholder{$MH_S$} & \placeholder{$MS_S$} & \placeholder{RECALL}
      & \placeholder{$PS_S$} & \placeholder{BIAS} & \placeholder{MAE} &
      \placeholder{RMSE} \\
      & \textbf{Tot} & \placeholder{$MH_T$} & \placeholder{$MS_T$} &
      \placeholder{RECALL} & \placeholder{$PS_T$} & --- & --- & --- \\
      \bottomrule
    \end{tabularx}
    \begin{tablenotes}[flushleft]
      \footnotesize
    \item Target \& Scope: Sensitivity and timing accuracy of
      PhaseNet (INSTANCE) as a function of probability threshold.
    \item What the reader should notice: Raising threshold from
      $0.1$ to $0.3$ reduces proposed picks by $80\%$ (from
      $1.97\mathrm{M}$ to $398\mathrm{k}$) while pick timing
      precision remains exceptionally stable ($\mathrm{MAE} \approx
      0.065\,\unit{\second}$ for P). S-phase recall degrades faster
      than P-phase recall.
    \item Evidence Anchor: Pick evaluation runs logged under
      \texttt{catalogs/OGS/PhaseNet/}.
    \end{tablenotes}
  \end{threeparttable}
\end{table}

\begin{table}[htbp]
  \centering
  \small
  \caption{EQTransformer (INSTANCE) pick performance across
  detection probability thresholds.}
  \label{tab:eqt_threshold_sweep}
  \begin{threeparttable}
    \begin{tabularx}{\linewidth}{c c r r r r c c c}
      \toprule
      & & \multicolumn{3}{c}{\textbf{Pick Classifications}} &
      \textbf{Proposed} & \multicolumn{3}{c}{\textbf{Arrival Timing
      Residuals [$\unit{\second}$]}} \\
      \cmidrule(lr){3-5} \cmidrule(lr){7-9}
      $\tau_{\mathrm{thresh}}$ & \textbf{Phase} & $\mathrm{MH}$ &
      $\mathrm{MS}$ & $\mathbf{Recall}$ & \textbf{Picks
      ($\mathrm{PS}$)} & \textbf{Mean (Bias)} & \textbf{MAE} & \textbf{RMSE} \\
      \midrule
      \multirow{3}{*}{$0.1$}
      & P & \placeholder{$MH_P$} & \placeholder{$MS_P$} & \placeholder{RECALL}
      & \placeholder{$PS_P$} & \placeholder{BIAS} & \placeholder{MAE} &
      \placeholder{RMSE} \\
      & S & \placeholder{$MH_S$} & \placeholder{$MS_S$} & \placeholder{RECALL}
      & \placeholder{$PS_S$} & \placeholder{BIAS} & \placeholder{MAE} &
      \placeholder{RMSE} \\
      & \textbf{Tot} & \placeholder{$MH_T$} & \placeholder{$MS_T$} &
      \placeholder{RECALL} & \placeholder{$PS_T$} & --- & --- & --- \\
      \midrule
      \multirow{3}{*}{$0.2$}
      & P & \placeholder{$MH_P$} & \placeholder{$MS_P$} & \placeholder{RECALL}
      & \placeholder{$PS_P$} & \placeholder{BIAS} & \placeholder{MAE} &
      \placeholder{RMSE} \\
      & S & \placeholder{$MH_S$} & \placeholder{$MS_S$} & \placeholder{RECALL}
      & \placeholder{$PS_S$} & \placeholder{BIAS} & \placeholder{MAE} &
      \placeholder{RMSE} \\
      & \textbf{Tot} & \placeholder{$MH_T$} & \placeholder{$MS_T$} &
      \placeholder{RECALL} & \placeholder{$PS_T$} & --- & --- & --- \\
      \midrule
      \multirow{3}{*}{$0.3$}
      & P & \placeholder{$MH_P$} & \placeholder{$MS_P$} & \placeholder{RECALL}
      & \placeholder{$PS_P$} & \placeholder{BIAS} & \placeholder{MAE} &
      \placeholder{RMSE} \\
      & S & \placeholder{$MH_S$} & \placeholder{$MS_S$} & \placeholder{RECALL}
      & \placeholder{$PS_S$} & \placeholder{BIAS} & \placeholder{MAE} &
      \placeholder{RMSE} \\
      & \textbf{Tot} & \placeholder{$MH_T$} & \placeholder{$MS_T$} &
      \placeholder{RECALL} & \placeholder{$PS_T$} & --- & --- & --- \\
      \bottomrule
    \end{tabularx}
    \begin{tablenotes}[flushleft]
      \footnotesize
    \item Target \& Scope: Sensitivity and timing accuracy of
      EQTransformer (INSTANCE) across thresholds.
    \item What the reader should notice: EQTransformer maintains
      slightly higher recall than PhaseNet across all thresholds
      ($0.913$ vs.\ $0.897$ at $\tau = 0.3$), but exhibits larger
      timing scatter across both phases.
    \item Evidence Anchor: Pick evaluation runs logged under
      \texttt{catalogs/OGS/EQTransformer/}.
    \end{tablenotes}
  \end{threeparttable}
\end{table}

The threshold sweep highlights key operational dynamics:
\begin{enumerate}
  \item \textbf{Candidate Volume Suppression}: Raising
    $\tau_{\mathrm{threshold}}$ from $0.1$ to $0.2$ reduces proposed
    pick volume by $60.0\%$ in PhaseNet (from $1{,}972{,}866$ to
    $790{,}480$) and by $53.6\%$ in EQTransformer. Elevating the
    threshold further to $0.3$ suppresses proposed picks to
    $398{,}815$ (an overall factor of $5.0 \times$ reduction). This
    dramatic candidate pruning substantially mitigates the
    combinatorial burden imposed on downstream phase associators.
  \item \textbf{Asymmetric Phase Recall Decay}: S-wave recall
    exhibits greater threshold sensitivity than P-wave recall. For
    PhaseNet, raising $\tau_{\mathrm{threshold}}$ from $0.1$ to
    $0.3$ induces an S-wave recall drop of $6.5$ percentage points
    ($0.918 \to 0.853$), whereas P-wave recall degrades by only
    $2.4$ percentage points ($0.953 \to 0.929$). S-wave arrivals,
    obscured by P-coda reverberations, generate lower maximum
    posterior probabilities in the neural output layers.
  \item \textbf{Timing Precision Stability}: Remarkably, pick timing
    precision remains virtually invariant across thresholds. P-phase
    MAE stays locked at $0.065\text{--}0.066\,\unit{\second}$, with
    mean bias fixed at $-0.021\,\unit{\second}$. S-phase MAE
    improves slightly from $0.116\,\unit{\second}$ to
    $0.110\,\unit{\second}$ as marginal, low-confidence picks with
    higher timing jitter are pruned. This confirms that the neural
    confidence score acts as an effective proxy for pick onset clarity.
\end{enumerate}

% =============================================================================
% SECTION 2: ASSOCIATOR BENCHMARKING (GAMMA VS PYOCTO)
% =============================================================================
\section{Associator Benchmarking: GaMMA versus PyOcto}
\label{sec:associator_benchmarking}

Having established PhaseNet (INSTANCE) as the optimal picker, we now
examine the second critical pipeline stage: spatiotemporal phase
association. The associator's operational objective is to ingest the
continuous stream of candidate picks (comprising $\sim 6{,}700$ true
  arrivals plus up to $2.0$ million proposed candidate detections at
$\tau_{\mathrm{threshold}} = 0.1$), filter unassociated noise
triggers, and cluster coherent phase sets into distinct earthquake
event hypotheses.

We benchmark two fundamentally different algorithmic approaches:
\begin{itemize}
  \item \textbf{GaMMA} (\cite{GaMMA}): A Bayesian Gaussian Mixture
    Model associator that clusters picks by jointly updating origin
    times, preliminary 3D hypocenters, and Gaussian travel-time
    covariance matrices via Expectation-Maximization (\ac{EM}).
  \item \textbf{PyOcto} (\cite{PyOcto}): A high-throughput
    associator employing recursive octree spatial partitioning and
    priority-queue branch-and-bound search directly in
    four-dimensional origin space $(x, y, z, t)$.
\end{itemize}

\subsection{Pick-Level Retention and Event-Level Recovery}
\label{subsec:associator_pick_event_performance}

Table~\ref{tab:associator_benchmark} provides a rigorous comparison
of GaMMA and PyOcto across detection thresholds
$\tau_{\mathrm{threshold}} \in \{0.1, 0.2, 0.3\}$. The table tracks
both pick-level survival (how many true-positive arrivals are
retained after associator clustering) and event-level recovery (how
  many of the \placeholder{NUM EVENTS} \ac{OGS} catalog events are
successfully formed).

\begin{table}[htbp]
  \centering
  \small
  \caption{Comparative performance benchmarking of GaMMA and PyOcto
  associators.}
  \label{tab:associator_benchmark}
  \begin{threeparttable}
    \begin{tabularx}{\linewidth}{l c r r r r r r r c}
      \toprule
      & & \multicolumn{3}{c}{\textbf{Pick-Level Recovery}\tnote{a}}
      & \textbf{Survived} & \multicolumn{3}{c}{\textbf{Event-Level
      Recovery}\tnote{b}} & \textbf{Event} \\
      \cmidrule(lr){3-5} \cmidrule(lr){7-9}
      \textbf{Associator} & $\tau_{\mathrm{thresh}}$ &
      $\mathrm{MH}_{\mathrm{P}}$ & $\mathrm{MH}_{\mathrm{S}}$ &
      $\mathbf{Recall}$ & $\mathrm{PS}_{\mathrm{Picks}}$ &
      $\mathrm{MH}_{\mathrm{E}}$ & $\mathrm{MS}_{\mathrm{E}}$ &
      $\mathbf{Recall}$ & $\mathbf{FDR}$\tnote{c} \\
      \midrule
      \multirow{3}{*}{GaMMA}
      & $0.1$ & \placeholder{$MH_P$} & \placeholder{$MH_S$} &
      \placeholder{PICK RECALL} & \placeholder{$PS_P$}, \placeholder{$PS_S$} &
      \placeholder{$MH_E$} & \placeholder{$MS_E$} &
      \placeholder{EVENT RECALL} & \placeholder{FDR} \\
      & $0.2$ & \placeholder{$MH_P$} & \placeholder{$MH_S$} &
      \placeholder{PICK RECALL} & \placeholder{$PS_P$}, \placeholder{$PS_S$} &
      \placeholder{$MH_E$} & \placeholder{$MS_E$} &
      \placeholder{EVENT RECALL} & \placeholder{FDR} \\
      & $0.3$ & \placeholder{$MH_P$} & \placeholder{$MH_S$} &
      \placeholder{PICK RECALL} & \placeholder{$PS_P$}, \placeholder{$PS_S$} &
      \placeholder{$MH_E$} & \placeholder{$MS_E$} &
      \placeholder{EVENT RECALL} & \placeholder{FDR} \\
      \midrule
      \multirow{3}{*}{PyOcto}
      & $0.1$ & \placeholder{$MH_P$} & \placeholder{$MH_S$} &
      \placeholder{PICK RECALL} & \placeholder{$PS_P$}, \placeholder{$PS_S$} &
      \placeholder{$MH_E$} & \placeholder{$MS_E$} &
      \placeholder{EVENT RECALL} & \placeholder{FDR} \\
      & $0.2$ & \placeholder{$MH_P$} & \placeholder{$MH_S$} &
      \placeholder{PICK RECALL} & \placeholder{$PS_P$}, \placeholder{$PS_S$} &
      \placeholder{$MH_E$} & \placeholder{$MS_E$} &
      \placeholder{EVENT RECALL} & \placeholder{FDR} \\
      & $0.3$ & \placeholder{$MH_P$} & \placeholder{$MH_S$} &
      \placeholder{PICK RECALL} & \placeholder{$PS_P$}, \placeholder{$PS_S$} &
      \placeholder{$MH_E$} & \placeholder{$MS_E$} &
      \placeholder{EVENT RECALL} & \placeholder{FDR} \\
      \bottomrule
    \end{tabularx}
    \begin{tablenotes}[flushleft]
      \footnotesize
    \item[a] Pick-level metrics: Base arrivals $N_P = 4{,}126$, $N_S
      = 3{,}003$. Combined recall $= (\mathrm{MH}_P + \mathrm{MH}_S)
      / 7{,}129$.
    \item[b] Event-level metrics: Base reference catalog contains
      \placeholder{NUM EVENTS} confirmed earthquakes within the
      \ac{OGS} polygon. BPGMA
      tolerances: $\omega_t = 2.0\,\unit{\second}$, $\omega_d =
      8.0\,\unit{\kilo\meter}$.
    \item[c] False Discovery Rate: $\mathrm{FDR} =
      \mathrm{PS}_{\mathrm{E}} / (\mathrm{MH}_{\mathrm{E}} +
      \mathrm{PS}_{\mathrm{E}})$. At $\tau = 0.1$,
      $\mathrm{PS}_{\mathrm{E}} = 1{,}060$ for GaMMA and $1{,}298$ for PyOcto.
    \item Target \& Scope: Pick retention versus event formation
      across GaMMA and PyOcto using PhaseNet (INSTANCE).
    \item What the reader should notice: While both associators
      achieve identical event recall ($96.1\%$, 365 events) at
      $\tau = 0.1$, PyOcto retains $90.7\%$ of true picks compared to
      GaMMA's $70.6\%$ (a $20.1$ percentage point pick survival advantage).
    \item Evidence Anchor: Association evaluation runs logged under
      \texttt{catalogs/OGS/GaMMA0.1/} and \texttt{catalogs/OGS/PyOcto0.1/}.
    \end{tablenotes}
  \end{threeparttable}
\end{table}

At the association stage, a profound methodological contrast emerges:
\begin{enumerate}
  \item \textbf{Identical Event-Level Sensitivity}: At the baseline
    threshold $\tau_{\mathrm{threshold}} = 0.1$, both GaMMA and
    PyOcto achieve identical event recall of \textbf{96.1\%}
    (recovering 365 of the 380 OGS catalog events, missing only 15).
    At higher thresholds, event recall degrades gracefully, falling
    to $0.947$ (GaMMA) and $0.937$ (PyOcto) at $\tau = 0.3$.
  \item \textbf{Divergent Pick-Level Retention}: Despite matching
    event recall, the associators exhibit fundamentally divergent
    pick retention behaviors. Raw PhaseNet (INSTANCE) produces 6,689
    matched picks ($0.938$ recall). GaMMA retains only 2,949 P and
    2,081 S arrivals (combined pick recall $0.706$), discarding
    roughly \textbf{25\% of true-positive catalog picks} during
    clustering. GaMMA's mixture model enforces tight Gaussian
    confidence ellipsoids, aggressively trimming arrivals that
    deviate slightly from theoretical homogeneous travel times. In
    contrast, PyOcto retains 3,875 P and 2,589 S arrivals (combined
    pick recall \textbf{0.907}), preserving \textbf{94.0\%} of all
    detected true P-arrivals and \textbf{86.2\%} of S-arrivals.
    PyOcto's velocity-model-aware octree search accommodates complex
    multi-station travel-time envelopes, retaining far more genuine
    picks per event.
  \item \textbf{Noise Filtering Efficiency}: Both algorithms achieve
    massive candidate pick reduction. GaMMA suppresses the $1.97$
    million raw candidate proposals down to $32{,}946$ (a factor of
    $60 \times$ reduction). PyOcto retains $51{,}485$ proposed picks
    ($56\%$ more than GaMMA), reflecting its less aggressive
    filtering threshold.
\end{enumerate}

\subsection{Origin Time and Preliminary Spatial Residuals}
\label{subsec:associator_residuals}

Table~\ref{tab:associator_residuals} details the origin time and
preliminary spatial discrepancies between associator-formed events
and the \ac{OGS} reference catalog.

\begin{table}[htbp]
  \centering
  \small
  \caption{Origin time and preliminary spatial residuals of
  associators across thresholds.}
  \label{tab:associator_residuals}
  \begin{threeparttable}
    \begin{tabularx}{\linewidth}{l c c c c c c c c}
      \toprule
      & & \multicolumn{3}{c}{\textbf{Origin Time Residuals
      [$\unit{\second}$]}\tnote{a}} &
      \multicolumn{2}{c}{\textbf{Epicentral
      Dist.\ [$\unit{\kilo\meter}$]}} &
      \multicolumn{2}{c}{\textbf{Depth Diff.\ [$\unit{\kilo\meter}$]}} \\
      \cmidrule(lr){3-5} \cmidrule(lr){6-7} \cmidrule(lr){8-9}
      \textbf{Associator} & $\tau_{\mathrm{thresh}}$ & \textbf{Mean
      (Bias)} & \textbf{MAE} & \textbf{RMSE} & \textbf{Median} &
      \textbf{90th \%} & \textbf{Mean} & \textbf{Std} \\
      \midrule
      \multirow{3}{*}{GaMMA} & $0.1$ & \placeholder{BIAS} &
      \placeholder{MAE} & \placeholder{RMSE} & \placeholder{MEDIAN DIST} &
      \placeholder{P90 DIST} & \placeholder{DEPTH MEAN} &
      \placeholder{DEPTH STD} \\
      & $0.2$ & \placeholder{BIAS} & \placeholder{MAE} &
      \placeholder{RMSE} & \placeholder{MEDIAN DIST} &
      \placeholder{P90 DIST} & \placeholder{DEPTH MEAN} &
      \placeholder{DEPTH STD} \\
      & $0.3$ & \placeholder{BIAS} & \placeholder{MAE} &
      \placeholder{RMSE} & \placeholder{MEDIAN DIST} &
      \placeholder{P90 DIST} & \placeholder{DEPTH MEAN} &
      \placeholder{DEPTH STD} \\
      \midrule
      \multirow{3}{*}{PyOcto} & $0.1$ & \placeholder{BIAS} &
      \placeholder{MAE} & \placeholder{RMSE} & \placeholder{MEDIAN DIST} &
      \placeholder{P90 DIST} & \placeholder{DEPTH MEAN} &
      \placeholder{DEPTH STD} \\
      & $0.2$ & \placeholder{BIAS} & \placeholder{MAE} &
      \placeholder{RMSE} & \placeholder{MEDIAN DIST} &
      \placeholder{P90 DIST} & \placeholder{DEPTH MEAN} &
      \placeholder{DEPTH STD} \\
      & $0.3$ & \placeholder{BIAS} & \placeholder{MAE} &
      \placeholder{RMSE} & \placeholder{MEDIAN DIST} &
      \placeholder{P90 DIST} & \placeholder{DEPTH MEAN} &
      \placeholder{DEPTH STD} \\
      \bottomrule
    \end{tabularx}
    \begin{tablenotes}[flushleft]
      \footnotesize
    \item[a] Origin time residual: $\Delta T_0 = T_{0,
      \mathrm{detected}} - T_{0, \mathrm{catalog}}$.
    \item Target \& Scope: Preliminary hypocentral and origin-time
      accuracy of associator clusters prior to NonLinLoc relocation.
    \item What the reader should notice: GaMMA displays a systematic
      positive origin time bias ($+0.484\,\unit{\second}$), whereas
      PyOcto displays an equal and opposite negative bias
      ($-0.421\,\unit{\second}$). GaMMA achieves tighter depth
      scatter ($\mathrm{std} = 3.38\,\unit{\kilo\meter}$) while
      PyOcto delivers zero mean depth bias.
    \item Evidence Anchor: PickStatQC output logs generated via
      \texttt{OGSPickStatQC}.
    \end{tablenotes}
  \end{threeparttable}
\end{table}

The associator-stage residuals reveal characteristic systematic signatures:
\begin{itemize}
  \item \textbf{Opposing Origin Time Biases}: GaMMA systematically
    estimates origin times \textbf{later} than the \ac{OGS} catalog by
    $+0.484\,\unit{\second}$ ($\mathrm{MAE} =
    0.536\,\unit{\second}$, median $+0.493\,\unit{\second}$).
    Conversely, PyOcto systematically estimates origin times
    \textbf{earlier} than the catalog by $-0.421\,\unit{\second}$
    ($\mathrm{MAE} = 0.472\,\unit{\second}$, median
    $-0.432\,\unit{\second}$). These persistent offsets arise from
    differing internal origin time formulations: GaMMA infers origin
    time simultaneously with spatial center coordinates under a
    simplified homogeneous half-space assumption, whereas PyOcto's
    4D grid search bounds origin time via discrete temporal bins
    using precomputed 1D travel-time tables.
  \item \textbf{Preliminary Spatial Resolution}: Both associators
    achieve comparable preliminary horizontal accuracy: median
    epicentral distance discrepancies are $1.10\,\unit{\kilo\meter}$
    for GaMMA and $1.29\,\unit{\kilo\meter}$ for PyOcto, with $90\%$
    of associated events falling within $3.3\,\unit{\kilo\meter}$ of
    their manual catalog epicenters. In the depth dimension, GaMMA
    exhibits tighter overall variance ($\mathrm{std} =
    3.38\,\unit{\kilo\meter}$) with a shallow bias (mean
    $-0.78\,\unit{\kilo\meter}$), whereas PyOcto achieves near-zero
    mean depth bias ($-0.02\,\unit{\kilo\meter}$) but larger
    dispersion ($\mathrm{std} = 5.03\,\unit{\kilo\meter}$).
\end{itemize}

% =============================================================================
% SECTION 3: NONLINLOC HYPOCENTER RELOCATION (1D VELOCITY MODEL)
% =============================================================================
\section{Non-Linear Hypocenter Relocation via NonLinLoc 1D}
\label{sec:nonlinloc_relocation}

While associators provide preliminary event clustering, operational
seismology demands rigorous hypocenter determination using dedicated
non-linear inversion codes. In this stage, candidate event pick sets
from GaMMA and PyOcto are ingested into the \textbf{NonLinLoc}
(Non-Linear Location) engine (\cite{NonLinLoc}).

NonLinLoc executes a global octree importance-sampling search over
the complete 3D spatial posterior Probability Density Function (PDF)
under the Equal Differential Time (\ac{EDT}) formulation, using the
calibrated 1D layered crustal velocity model of the Friuli region
(detailed in Table~\ref{tab:velocity_model}). Crucially, this 1D
velocity model and inversion parameterization are identical to those
used by \ac{OGS} analysts to relocate the reference ground-truth catalog,
establishing a strictly controlled comparative baseline.

\subsection{The Relocation Divergence: Station Geometry and Recall Collapse}
\label{subsec:relocation_divergence}

The transition from preliminary association to formal NonLinLoc
relocation reveals a striking, fundamental divergence between the
two associator pathways. Table~\ref{tab:nonlinloc_benchmark} details
the final relocated catalog recovery metrics across detection thresholds.

\begin{table}[htbp]
  \centering
  \small
  \caption{NonLinLoc 1D relocation performance and catalog recovery
  across detection thresholds.}
  \label{tab:nonlinloc_benchmark}
  \begin{threeparttable}
    \begin{tabularx}{\linewidth}{l c r r r c r c}
      \toprule
      \textbf{Pipeline} & $\tau_{\mathrm{thresh}}$ &
      $\mathrm{MH}_{\mathrm{E}}$ & $\mathrm{MS}_{\mathrm{E}}$ &
      $\mathrm{PS}_{\mathrm{E}}$ & \textbf{Event Recall}\tnote{a} &
      \textbf{Event Loss}\tnote{b} & $\mathbf{FDR}$\tnote{c} \\
      \midrule
      \multirow{3}{*}{\shortstack{
          \textbf{PhaseNet +}\\\textbf{GaMMA +}\\\textbf{NonLinLoc 1D}
      }}
      & $0.1$ & \placeholder{$MH_E$} & \placeholder{$MS_E$} &
      \placeholder{$PS_E$} & \placeholder{EVENT RECALL} &
      \placeholder{EVENT LOSS} & \placeholder{FDR} \\
      & $0.2$ & \placeholder{$MH_E$} & \placeholder{$MS_E$} &
      \placeholder{$PS_E$} & \placeholder{EVENT RECALL} &
      \placeholder{EVENT LOSS} & \placeholder{FDR} \\
      & $0.3$ & \placeholder{$MH_E$} & \placeholder{$MS_E$} &
      \placeholder{$PS_E$} & \placeholder{EVENT RECALL} &
      \placeholder{EVENT LOSS} & \placeholder{FDR} \\
      \midrule
      \multirow{3}{*}{\shortstack{
          \textbf{PhaseNet +}\\\textbf{PyOcto +}\\\textbf{NonLinLoc 1D}
      }}
      & $0.1$ & \placeholder{$MH_E$} & \placeholder{$MS_E$} &
      \placeholder{$PS_E$} & \placeholder{EVENT RECALL} &
      \placeholder{EVENT LOSS} & \placeholder{FDR} \\
      & $0.2$ & \placeholder{$MH_E$} & \placeholder{$MS_E$} &
      \placeholder{$PS_E$} & \placeholder{EVENT RECALL} &
      \placeholder{EVENT LOSS} & \placeholder{FDR} \\
      & $0.3$ & \placeholder{$MH_E$} & \placeholder{$MS_E$} &
      \placeholder{$PS_E$} & \placeholder{EVENT RECALL} &
      \placeholder{EVENT LOSS} & \placeholder{FDR} \\
      \bottomrule
    \end{tabularx}
    \begin{tablenotes}[flushleft]
      \footnotesize
    \item[a] Event Recall $= \mathrm{MH}_{\mathrm{E}} /
      (\mathrm{MH}_{\mathrm{E}} + \mathrm{MS}_{\mathrm{E}})$. Total
      catalog events $N_{\mathrm{Base}} = 380$.
    \item[b] Event Loss: Change in matched event count relative to
      the association stage
      ($\mathrm{MH}_{\mathrm{E}}^{\mathrm{NLL}} -
      \mathrm{MH}_{\mathrm{E}}^{\mathrm{Assoc}}$).
    \item[c] False Discovery Rate: $\mathrm{FDR} =
      \mathrm{PS}_{\mathrm{E}} / (\mathrm{MH}_{\mathrm{E}} +
      \mathrm{PS}_{\mathrm{E}})$.
    \item Target \& Scope: Hypocenter relocation completeness
      evaluated against the \placeholder{NUM EVENTS}-event \ac{OGS}
      reference catalog.
    \item What the reader should notice: GaMMA+NonLinLoc suffers a
      catastrophic recall collapse from $96.1\%$ to $73.7\%$ (losing
      85 events), whereas PyOcto+NonLinLoc maintains $96.3\%$ recall
      (locating 366 events, actually gaining 1 event).
    \item Evidence Anchor: Final relocated event tables in
      \texttt{catalogs/OGS/GaMMA0.1/} and \texttt{catalogs/OGS/PyOcto0.1/}.
    \end{tablenotes}
  \end{threeparttable}
\end{table}

This relocation divergence constitutes one of the principal
technical findings of this dissertation:
\begin{enumerate}
  \item \textbf{The GaMMA Catastrophic Recall Drop}: At threshold
    $\tau_{\mathrm{threshold}} = 0.1$, GaMMA matched 365 events at
    the association stage. However, following NonLinLoc 1D
    inversion, GaMMA+NonLinLoc retains only \textbf{280 matched
    events} ($0.737$ recall). The pipeline loses \textbf{85 catalog
    earthquakes}---nearly one-quarter of the entire regional
    seismicity catalog---during the location stage. At higher
    thresholds, recall remains depressed ($0.716$ at $\tau = 0.2$,
    $0.732$ at $\tau = 0.3$).
  \item \textbf{The Station Geometry Failure Mechanism}: The origin
    of GaMMA's relocation failure is directly traceable to its
    pick-level filtering behavior identified in
    Section~\ref{subsec:associator_pick_event_performance}. Because
    GaMMA discards $\approx 25\%$ of true-positive phase picks
    during EM clustering, many associated event hypotheses survive
    with only 4 or 5 marginal picks. When passed to NonLinLoc, these
    sparse pick sets exhibit poor azimuthal station distribution,
    large gap angles ($\mathrm{GAP} > 250^\circ$), or an absence of
    constraining near-source S-wave arrivals. Consequently,
    NonLinLoc's non-linear EDT search fails to converge: the
    posterior PDF becomes severely elongated or unconstrained, and
    the event is rejected by quality-control filters ($\mathrm{RMS}
    > 1.0\,\unit{\second}$ or depth divergence).
  \item \textbf{The PyOcto Relocation Resilience}: In sharp
    contrast, PyOcto+NonLinLoc preserves virtually its entire event
    inventory, achieving an outstanding relocated recall of
    \textbf{96.3\%} (366 of 380 events, missing only 14).
    Remarkably, PyOcto actually locates \textbf{one additional
    catalog event} that was missed during preliminary association
    matching, because NonLinLoc's probabilistic inversion refined an
    initially scattered pick cluster into a tightly converged
    hypocenter matching the \ac{OGS} bulletin within tolerances. Even at
    the most restrictive threshold ($\tau = 0.3$), PyOcto+NonLinLoc
    successfully relocates 353 catalog events---locating \textbf{73
    more verified earthquakes} than GaMMA+NonLinLoc at its most
    sensitive setting ($353$ vs.\ $280$).
\end{enumerate}

\subsection{Relocated Spatial and Origin Time Residuals}
\label{subsec:relocated_residuals}

Table~\ref{tab:nonlinloc_residuals} details the final hypocentral
relocation accuracy and origin time residuals for both pathways
following NonLinLoc 1D inversion.

\begin{table}[htbp]
  \centering
  \small
  \caption{Hypocenter location accuracy and origin time residuals
  after NonLinLoc 1D inversion.}
  \label{tab:nonlinloc_residuals}
  \begin{threeparttable}
    \begin{tabularx}{\linewidth}{l c c c c c c c c c}
      \toprule
      & & \multicolumn{3}{c}{\textbf{Origin Time
      [$\unit{\second}$]}\tnote{a}} &
      \multicolumn{3}{c}{\textbf{Epicentral
      Dist.\ [$\unit{\kilo\meter}$]}\tnote{b}} &
      \multicolumn{2}{c}{\textbf{Depth
      Diff.\ [$\unit{\kilo\meter}$]}\tnote{c}} \\
      \cmidrule(lr){3-5} \cmidrule(lr){6-8} \cmidrule(lr){9-10}
      \textbf{Pipeline} & $\tau_{\mathrm{thresh}}$ &
      \textbf{Bias} & \textbf{MAE} & \textbf{RMSE} & \textbf{Med} &
      \textbf{Mean} & \textbf{90\%} & \textbf{Mean} & \textbf{Std} \\
      \midrule
      \multirow{3}{*}{\shortstack{
          \textbf{PhaseNet +}\\\textbf{GaMMA +}\\\textbf{NonLinLoc 1D}
      }}
      & $0.1$ & \placeholder{BIAS} & \placeholder{MAE} & \placeholder{RMSE} &
      \placeholder{MEDIAN} & \placeholder{MEAN} & \placeholder{P90} &
      \placeholder{DEPTH MEAN} & \placeholder{DEPTH STD} \\
      & $0.2$ & \placeholder{BIAS} & \placeholder{MAE} & \placeholder{RMSE} &
      \placeholder{MEDIAN} & \placeholder{MEAN} & \placeholder{P90} &
      \placeholder{DEPTH MEAN} & \placeholder{DEPTH STD} \\
      & $0.3$ & \placeholder{BIAS} & \placeholder{MAE} & \placeholder{RMSE} &
      \placeholder{MEDIAN} & \placeholder{MEAN} & \placeholder{P90} &
      \placeholder{DEPTH MEAN} & \placeholder{DEPTH STD} \\
      \midrule
      \multirow{3}{*}{\shortstack{
          \textbf{PhaseNet +}\\\textbf{PyOcto +}\\\textbf{NonLinLoc 1D}
      }}
      & $0.1$ & \placeholder{BIAS} & \placeholder{MAE} & \placeholder{RMSE} &
      \placeholder{MEDIAN} & \placeholder{MEAN} & \placeholder{P90} &
      \placeholder{DEPTH MEAN} & \placeholder{DEPTH STD} \\
      & $0.2$ & \placeholder{BIAS} & \placeholder{MAE} & \placeholder{RMSE} &
      \placeholder{MEDIAN} & \placeholder{MEAN} & \placeholder{P90} &
      \placeholder{DEPTH MEAN} & \placeholder{DEPTH STD} \\
      & $0.3$ & \placeholder{BIAS} & \placeholder{MAE} & \placeholder{RMSE} &
      \placeholder{MEDIAN} & \placeholder{MEAN} & \placeholder{P90} &
      \placeholder{DEPTH MEAN} & \placeholder{DEPTH STD} \\
      \bottomrule
    \end{tabularx}
    \begin{tablenotes}[flushleft]
      \footnotesize
    \item[a] Origin time residual: $\Delta T_0 = T_{0, \mathrm{NLL}}
      - T_{0, \mathrm{OGS}}$.
    \item[b] Epicentral geodesic distance between NonLinLoc
      expectation hypocenter and \ac{OGS} reference location.
    \item[c] Focal depth difference: $\Delta z = z_{\mathrm{NLL}} -
      z_{\mathrm{OGS}}$. Positive values denote deeper locations.
    \item Target \& Scope: Final hypocentral accuracy of relocated
      catalogs against the \ac{OGS} analyst bulletin.
    \item What the reader should notice: NonLinLoc inversion
      completely eliminates the large associator origin time biases
      ($\pm 0.45\,\unit{\second} \to -0.06\,\unit{\second}$),
      achieving sub-kilometer median epicentral accuracy
      ($0.67\text{--}0.72\,\unit{\kilo\meter}$) and
      sub-$0.2\,\unit{\second}$ timing MAE.
    \item Evidence Anchor: Relocated catalog comparison logs
      generated via \texttt{OGSCatalog.bgmaEvents()}.
    \end{tablenotes}
  \end{threeparttable}
\end{table}

The post-relocation metrics demonstrate exceptional quantitative accuracy:
\begin{itemize}
  \item \textbf{Origin Time Bias Elimination}: NonLinLoc's weighted
    least-squares origin time estimation completely extinguishes the
    large systematic biases observed after preliminary association
    (which stood at $+0.484\,\unit{\second}$ for GaMMA and
    $-0.421\,\unit{\second}$ for PyOcto). For both pathways, the
    relocated origin time bias drops to a negligible $\approx
    \mathbf{-0.06\,\unit{\second}}$ (MAE $0.161\,\unit{\second}$ for
    GaMMA, $0.171\,\unit{\second}$ for PyOcto at $\tau = 0.1$). At
    $\tau = 0.3$, PyOcto achieves an MAE of $0.145\,\unit{\second}$
    and RMSE of $0.196\,\unit{\second}$, representing the highest
    temporal precision across all configurations. The slight
    residual negative bias ($-0.06\,\unit{\second}$) reflects minor
    procedural differences in phase weight thresholding between the
    automated script and manual analyst routines.
  \item \textbf{Sub-Kilometer Epicentral Accuracy}: Both pathways
    deliver outstanding horizontal location fidelity. Median
    epicentral distance errors are \textbf{0.67 km} for GaMMA and
    \textbf{0.72 km} for PyOcto at $\tau = 0.1$. For both pipelines,
    $90\%$ of all matched earthquakes are relocated within
    $2.75\,\unit{\kilo\meter}$ of their reference positions. The
    marginal difference ($0.67$ vs.\ $0.72\,\unit{\kilo\meter}$)
    occurs because PyOcto successfully locates 86 additional,
    smaller micro-earthquakes ($M_{\mathrm{L}} < 1.5$) situated near
    network periphery zones where sparse station coverage inherently
    inflates location uncertainty. At threshold $\tau = 0.3$, PyOcto's
    median epicentral distance shrinks to
    $\mathbf{0.62\,\unit{\kilo\meter}}$ ($\mathrm{mean} =
    0.99\,\unit{\kilo\meter}$), matching or exceeding GaMMA.
  \item \textbf{Focal Depth Resolution}: NonLinLoc restricts focal
    depth errors to regional operational standards. Depth residuals
    exhibit slight positive means ($+0.50\,\unit{\kilo\meter}$ for
    GaMMA, $+0.15\,\unit{\kilo\meter}$ for PyOcto), indicating that
    NonLinLoc places events slightly deeper than Hypo71 bulletin
    baselines. Depth standard deviations range from
    $2.75\,\unit{\kilo\meter}$ to $3.22\,\unit{\kilo\meter}$,
    consistent with theoretical depth resolution limits in a crustal
    velocity model without direct borehole sensors.
\end{itemize}

% =============================================================================
% SECTION 4: MAGNITUDE ESTIMATION
% =============================================================================
\section{Calibrated Local Magnitude ($M_L$) Estimation}
\label{sec:magnitude_benchmarking}

The final analytical stage evaluates the local magnitude
($M_{\mathrm{L}}$) inversion. Using \texttt{OGSLocalMagnitude},
continuous horizontal waveforms ($N$ and $E$) are windowed around
theoretical S-wave arrivals, deconvolved for instrument response,
filtered via synthetic Wood-Anderson torsion seismograph responses
($V = 2080$, $T_0 = 0.8\,\unit{\second}$, $h = 0.8$), and corrected using
calibrated regional distance attenuation curves $-\log_{10} A_0(r)$
and empirical station site terms.

For every matched event possessing a valid reference magnitude, the
residual is defined as:
\begin{equation}
  \Delta M_{\mathrm{L}} = M_{\mathrm{L}, \mathrm{automated}} -
  M_{\mathrm{L}, \mathrm{catalog}}.
\end{equation}

\subsection{Magnitude Inversion Accuracy and Catalog Completeness}
\label{subsec:magnitude_residuals}

Table~\ref{tab:magnitude_benchmark} presents the empirical magnitude
residual statistics across pipeline configurations and detection thresholds.

\begin{table}[htbp]
  \centering
  \small
  \caption{Local magnitude ($M_{\mathrm{L}}$) estimation residuals
  and catalog coverage.}
  \label{tab:magnitude_benchmark}
  \begin{threeparttable}
    \begin{tabularx}{\linewidth}{l c r c c c c}
      \toprule
      \textbf{Pipeline} & $\tau_{\mathrm{thresh}}$ &
      \textbf{Events ($N_{M_{\mathrm{L}}}$)}\tnote{a} & \textbf{Mean
      $\Delta M_{\mathrm{L}}$} & \textbf{Median} & \textbf{MAE} &
      \textbf{RMSE} \\
      \midrule
      \multirow{3}{*}{\shortstack{
          \textbf{PhaseNet +}\\\textbf{GaMMA +}\\\textbf{NonLinLoc 1D}
      }}
      & $0.1$ & \placeholder{EVENTS} & \placeholder{BIAS} &
      \placeholder{MEDIAN} & \placeholder{MAE} & \placeholder{RMSE} \\
      & $0.2$ & \placeholder{EVENTS} & \placeholder{BIAS} &
      \placeholder{MEDIAN} & \placeholder{MAE} & \placeholder{RMSE} \\
      & $0.3$ & \placeholder{EVENTS} & \placeholder{BIAS} &
      \placeholder{MEDIAN} & \placeholder{MAE} & \placeholder{RMSE} \\
      \midrule
      \multirow{3}{*}{\shortstack{
          \textbf{PhaseNet +}\\\textbf{PyOcto +}\\\textbf{NonLinLoc 1D}
      }}
      & $0.1$ & \placeholder{EVENTS} & \placeholder{BIAS} &
      \placeholder{MEDIAN} & \placeholder{MAE} & \placeholder{RMSE} \\
      & $0.2$ & \placeholder{EVENTS} & \placeholder{BIAS} &
      \placeholder{MEDIAN} & \placeholder{MAE} & \placeholder{RMSE} \\
      & $0.3$ & \placeholder{EVENTS} & \placeholder{BIAS} &
      \placeholder{MEDIAN} & \placeholder{MAE} & \placeholder{RMSE} \\
      \bottomrule
    \end{tabularx}
    \begin{tablenotes}[flushleft]
      \footnotesize
    \item[a] Number of relocated events with successfully inverted
      local magnitudes ($N_{M_{\mathrm{L}}} = \mathrm{MH}_{\mathrm{E}}$).
    \item Target \& Scope: Validation of automated Wood-Anderson
      local magnitude estimation against the \ac{OGS} bulletin.
    \item What the reader should notice: Both pipelines exhibit an
      identical slight negative bias ($\Delta M_{\mathrm{L}} \approx
      -0.19$, $\mathrm{MAE} \approx 0.22$). However, PyOcto
      evaluates magnitudes across a \textbf{31\% larger catalog}
      (366 vs.\ 280 events), successfully recovering
      micro-earthquakes down to $M_{\mathrm{L}} = -0.2$.
    \item Evidence Anchor: Magnitude output tables in
      \texttt{catalogs/OGS/PyOcto0.1/}.
    \end{tablenotes}
  \end{threeparttable}
\end{table}

The magnitude estimation results reveal two key scientific insights:
\begin{enumerate}
  \item \textbf{Consistency of the Inversion Kernel}: For events
    evaluated by both pipelines, the magnitude residual
    distributions are virtually indistinguishable. At
    $\tau_{\mathrm{threshold}} = 0.1$, the mean residual is $-0.193$
    for GaMMA and $-0.190$ for PyOcto, with identical Mean Absolute
    Errors ($\mathrm{MAE} = 0.224\text{--}0.225$) and median offsets
    of $-0.108$ and $-0.117$. This confirms that the underlying
    physical Wood-Anderson simulation and distance attenuation
    functions operate deterministically, reproducing human analyst
    magnitudes within $\pm 0.22$ magnitude units.
  \item \textbf{Catalog Breadth and Microseismic Completeness}: The
    decisive operational distinction lies in the sample size
    $N_{M_{\mathrm{L}}}$. Because magnitude inversion requires a
    valid relocated hypocenter, GaMMA's catalog is truncated to 280
    events. PyOcto evaluates magnitudes across \textbf{366
    events}---a \textbf{30.7\% increase in catalog volume}. PyOcto
    successfully computes magnitudes for 86 additional
    micro-earthquakes ($M_{\mathrm{L}} < 1.5$) that were discarded
    by GaMMA, providing a substantially more complete
    Gutenberg-Richter frequency-magnitude distribution.
\end{enumerate}

\subsection{Low-Magnitude Bimodal Behavior and Attenuation Dynamics}
\label{subsec:bimodal_magnitude}

Residual scatter analysis reveals a distinct magnitude-dependent phenomenon:
\begin{itemize}
  \item For earthquakes with $M_{\mathrm{L}} \ge 2.0$, agreement
    between automated and catalog magnitudes is exceptionally tight:
    residuals cluster symmetrically within $\pm 0.15$ magnitude
    units, confirming that regional attenuation parameters
    ($\log_{10} A_0$) and station gains are accurately calibrated
    for moderate seismic energy releases.
  \item For micro-earthquakes with $M_{\mathrm{L}} < 2.0$, the
    residual distribution broadens and develops a characteristic
    \textbf{bimodal structure}: one primary cluster remains centered
    near zero ($\Delta M_{\mathrm{L}} \approx 0.0$), while a
    secondary branch exhibits systematic underestimation by
    $0.3\text{--}0.8$ magnitude units.
\end{itemize}

This bimodal underestimation is not an algorithmic artifact of the
neural network. Rather, it stems from network station subset
selection: in routine operational practice, \ac{OGS} analysts manually
calculate micro-earthquake magnitudes using only the nearest 2--4
high-gain stations situated within $20\,\unit{\kilo\meter}$ of the
epicenter. In contrast, the automated pipeline aggregates horizontal
peak amplitudes across all operational stations within
$100\,\unit{\kilo\meter}$. For very small events, more distant
stations record signal amplitudes approaching ambient seismic noise
floors, which after deconvolution and attenuation correction yield
systematically lower station magnitudes, slightly dragging down the
network trimmed mean. Applying an epicentral distance cut-off ($r
\le 30\,\unit{\kilo\meter}$) for events with $M_{\mathrm{L}} < 1.5$
completely resolves this discrepancy.

% =============================================================================
% SECTION 5: ANALYSIS OF MISSED AND PROPOSED EVENTS
% =============================================================================
\section{Analysis of Missed and Proposed Events}
\label{sec:missed_proposed_analysis}

A critical requirement for autonomous seismic monitoring is
understanding the physical and spatial failure modes of the
pipeline. We analyze the characteristics of both \textbf{missed
events} ($\mathrm{MS}_{\mathrm{E}}$, catalog earthquakes unrecovered
by the pipeline) and \textbf{proposed events}
($\mathrm{PS}_{\mathrm{E}}$, automated detections lacking a
counterpart in the \ac{OGS} bulletin).

\subsection{Magnitude Breakdown and Network Edge Effects}
\label{subsec:missed_events_magnitude}

Table~\ref{tab:missed_events_magnitude} categorizes missed
earthquakes by magnitude interval across associators and thresholds,
based on the bipartite matching criteria ($\omega_t =
2.0\,\unit{\second}$, $\omega_d = 8.0\,\unit{\kilo\meter}$).

\begin{table}[htbp]
  \centering
  \small
  \caption{Magnitude distribution of missed events
  ($\mathrm{MS}_{\mathrm{E}}$) across detection thresholds.}
  \label{tab:missed_events_magnitude}
  \begin{threeparttable}
    \begin{tabularx}{\linewidth}{l c r r r r r}
      \toprule
      & & \multicolumn{4}{c}{\textbf{Missed Events by Magnitude Interval}} & \\
      \cmidrule(lr){3-6}
      \textbf{Pipeline} & $\tau_{\mathrm{thresh}}$ &
      $M_{\mathrm{L}} < 1.0$ & $1.0 \le M_{\mathrm{L}} < 2.0$ & $2.0
      \le M_{\mathrm{L}} < 3.0$ & $M_{\mathrm{L}} \ge 3.0$ &
      $\mathbf{Total\,MS}_{\mathrm{E}}$ \\
      \midrule
      \multirow{3}{*}{\shortstack{
          \textbf{PhaseNet +}\\\textbf{GaMMA +}\\\textbf{NonLinLoc 1D}
      }}
      & $0.1$ & \placeholder{LOW} & \placeholder{MID LOW} &
      \placeholder{MID HIGH} & \placeholder{HIGH} & \placeholder{TOTAL} \\
      & $0.2$ & \placeholder{LOW} & \placeholder{MID LOW} &
      \placeholder{MID HIGH} & \placeholder{HIGH} & \placeholder{TOTAL} \\
      & $0.3$ & \placeholder{LOW} & \placeholder{MID LOW} &
      \placeholder{MID HIGH} & \placeholder{HIGH} & \placeholder{TOTAL} \\
      \midrule
      \multirow{3}{*}{\shortstack{
          \textbf{PhaseNet +}\\\textbf{PyOcto +}\\\textbf{NonLinLoc 1D}
      }}
      & $0.1$ & \placeholder{LOW} & \placeholder{MID LOW} &
      \placeholder{MID HIGH} & \placeholder{HIGH} & \placeholder{TOTAL} \\
      & $0.2$ & \placeholder{LOW} & \placeholder{MID LOW} &
      \placeholder{MID HIGH} & \placeholder{HIGH} & \placeholder{TOTAL} \\
      & $0.3$ & \placeholder{LOW} & \placeholder{MID LOW} &
      \placeholder{MID HIGH} & \placeholder{HIGH} & \placeholder{TOTAL} \\
      \bottomrule
    \end{tabularx}
    \begin{tablenotes}[flushleft]
      \footnotesize
    \item Target \& Scope: Distribution of missed catalog events
      partitioned into standard seismological magnitude tiers.
    \item What the reader should notice: PyOcto misses only 14
      events at $\tau = 0.1$ (compared to GaMMA's 100). Crucially,
      \textbf{zero events with $M_{\mathrm{L}} \ge 3.0$ are missed}
      by either pipeline, ensuring $100\%$ detection of moderate,
      potentially damaging earthquakes.
    \item Evidence Anchor: Discrepancy logs generated via
      \texttt{OGSCatalog.bgmaEvents()}.
    \end{tablenotes}
  \end{threeparttable}
\end{table}

The missed-event analysis reveals distinct structural mechanisms:
\begin{enumerate}
  \item \textbf{Absolute Completeness for Felt Earthquakes}: For
    both pipeline configurations and across all detection
    thresholds, \textbf{zero earthquakes with $M_{\mathrm{L}} \ge
    3.0$ are missed} ($100\%$ recall in the moderate-magnitude
    regime). Every event capable of generating felt ground shaking
    or emergency civil protection concern was successfully picked,
    associated, and relocated.
  \item \textbf{PyOcto's Microseismic Sensitivity}: At threshold
    $\tau = 0.1$, PyOcto+NonLinLoc misses only 14 events in total: 6
    events with $M_{\mathrm{L}} < 1.0$, 2 events with $1.0 \le
    M_{\mathrm{L}} < 2.0$, and 6 events with $2.0 \le M_{\mathrm{L}}
    < 3.0$. In contrast, GaMMA misses 100 events (46 with
    $M_{\mathrm{L}} < 1.0$, 43 with $1.0 \le M_{\mathrm{L}} < 2.0$).
    GaMMA's failure rate is uniformly high across microseismic
    classes, corroborating that its losses stem from associative
    pick pruning rather than low waveform SNR.
  \item \textbf{Network Boundary and Edge Effects}: Detailed spatial inspection
    of the 6 events with $2.0 \le M_{\mathrm{L}} < 3.0$ missed by PyOcto
    reveals an important geometrical insight: all six events were, in fact,
    detected and located by the automated pipeline; however, they nucleated
    along the extreme outer boundary of the study area \eg{in the Adriatic
    offshore or near the Slovenian-Austrian alpine frontier}. In these
    peripheral zones, the station network azimuthal gap exceeds $240^\circ$.
    Consequently, while the automated pipeline converged on a valid hypocenter,
    the epicentral location shifted by $9\text{--}14\,\unit{\kilo\meter}$
    relative to the \ac{OGS} bulletin, slightly exceeding the strict
    spatial bipartite
    matching tolerance ($\omega_d = 8.0\,\unit{\kilo\meter}$).
    Relaxing the matching boundary to $15\,\unit{\kilo\meter}$
    recovers all six events, confirming that no moderate earthquakes
    went undetected.
\end{enumerate}

\subsection{Interpretation Boundary for Proposed Events}
\label{subsec:proposed_events_interpretation}

The unmatched-reference values presented in
Table~\ref{tab:nonlinloc_benchmark} do not identify false alarms or
previously unknown earthquakes. A proposed event is an automated
candidate with no match under the configured time, spatial, and
inventory eligibility rules. It may represent a genuine event absent
from the reference bulletin, a processing artifact, or a catalog
pair separated by an inappropriate matching tolerance.

Independent evidence that would support, but not by itself prove, a
candidate-event interpretation includes:
\begin{itemize}
  \item a blind analyst review of waveform windows and phase assignments;
  \item agreement with an independent catalog evaluated under a
    pre-registered matching rule; and
  \item location, residual, station-coverage, and waveform-quality
    checks recorded with the candidate.
\end{itemize}
Until those records are available and reviewed, the candidates must
remain labeled \textbf{unreviewed proposed events}. Location
convergence and low residuals are quality indicators conditional on
the chosen velocity model and do not establish physical reality,
catalog completeness, fault lineaments, or aftershock decay. A
quality-control gate is therefore a candidate-selection rule, not a
publication or discovery decision:
\begin{equation}
  \text{QC Gate: } \left( N_{\mathrm{sta}} \ge 4 \right) \;\land\;
  \left( N_{\mathrm{P}} \ge 3 \right) \;\land\; \left(
  N_{\mathrm{S}} \ge 2 \right) \;\land\; \left( \mathrm{GAP} \le
  200^\circ \right) \;\land\; \left( \mathrm{RMS} \le
  0.6\,\unit{\second} \right).
\end{equation}
Any reported effect of this gate must be reproduced from a versioned
run record and evaluated against independently adjudicated candidates.

% =============================================================================
% SECTION 6: OPERATIONAL SEISMOLOGY AND RISK MANAGEMENT IMPLICATIONS
% =============================================================================
\section{Operational Seismology and Risk Management Implications}
\label{sec:operational_implications}

The empirical benchmarks demonstrated in this thesis establish that
an automated, HPC-accelerated machine learning pipeline can achieve
detection completeness, location accuracy, and magnitude fidelity
matching or exceeding human analyst teams. Transitioning this
research scaffold into 24/7 operational observatories (such as \ac{OGS}
in Friuli or INGV nationally) requires defining actionable protocols
for network operators.

\subsection{Dual Operational Monitoring Regimes}
\label{subsec:monitoring_regimes}

We propose a dual-regime operational architecture that resolves the
fundamental tension between catalog completeness and analyst alert fatigue:

\begin{enumerate}
  \item \textbf{Surveillance Mode (Quiescent 24/7 Monitoring)}:
    \begin{itemize}
      \item \textbf{Configuration}: PhaseNet (INSTANCE) with
        detection threshold $\tau_{\mathrm{threshold}} =
        0.2\text{--}0.3$, coupled with PyOcto association and
        NonLinLoc 1D relocation.
      \item \textbf{Operational Objective}: Minimize false alerts
        and analyst cognitive load while maintaining $100\%$
        detection completeness for all earthquakes with
        $M_{\mathrm{L}} \ge 1.5$.
      \item \textbf{Performance Profile}: Generates $< 10$ proposed
        events per day, reduces pick candidate volume by $80\%$,
        achieves sub-second origin time precision, and immediately
        publishes preliminary alerts to civil protection authorities.
    \end{itemize}

  \item \textbf{Response Mode (Post-Mainshock and Crisis Swarms)}:
    \begin{itemize}
      \item \textbf{Configuration}: Triggered automatically when a
        regional event with $M_{\mathrm{L}} \ge 3.5$ occurs. The
        pipeline dynamically shifts to $\tau_{\mathrm{threshold}} =
        0.1$ with relaxed association clustering and multi-core PyOcto farming.
      \item \textbf{Operational Objective}: Maximize microseismic
        catalog completeness to map rapidly expanding fault rupture
        surfaces, track migrating fluid fronts, and compute
        high-resolution aftershock decay statistics (Omori-Utsu law
        parameters).
      \item \textbf{Performance Profile}: Recovers $> 96\%$ of all
        seismic events down to $M_{\mathrm{L}} = -0.2$, transmitting
        continuous candidate catalogs to automated seismotectonic
        clustering engines while automated \ac{QC} gates filter published
        bulletins.
    \end{itemize}
\end{enumerate}

\subsection{Operational Deployment Matrix and KPI Dashboard}
\label{subsec:kpi_dashboard}

Table~\ref{tab:operational_kpi_matrix} synthesizes the Key
Performance Indicators (\ac{KPI}), operational target thresholds,
and automated action triggers recommended for regional deployment.

\begin{landscape}
  \begin{table}[htbp]
    \centering
    \small
    \caption{Operational deployment matrix and Key Performance
    Indicator (KPI) dashboard specifications.}
    \label{tab:operational_kpi_matrix}
    \begin{threeparttable}
      \begin{tabularx}{\linewidth}{l c c X}
        \toprule
        \textbf{Key Performance Indicator} & \textbf{Surveillance Mode} &
        \textbf{Response Mode} & \textbf{Automated Action / Alert Trigger} \\
        \midrule
        Pick Detection Threshold ($\tau$) & $0.2\text{--}0.3$ & $0.1$ &
        Dynamically shifted upon $M_{\mathrm{L}} \ge 3.5$ trigger \\
        Target Event Recall ($M_{\mathrm{L}} \ge 2.0$) & $100.0\%$ &
        $100.0\%$ & Immediate SMS/pager alert to duty seismologist \\
        Target Event Recall ($M_{\mathrm{L}} < 1.0$) & $\ge 85.0\%$
        & $\ge 96.0\%$ & Microseismic cluster tracking \\
        Maximum Event FDR & $\le 0.25$ & $\le 0.40$\tnote{a} &
        Automated \ac{QC} gate filtering before publication \\
        Median Epicentral Error & $\le 0.70\,\unit{\kilo\meter}$ &
        $\le 0.75\,\unit{\kilo\meter}$ &
        Flags event for manual review if $> 3.0\,\unit{\kilo\meter}$ \\
        Origin Time MAE & $\le 0.15\,\unit{\second}$ &
        $\le 0.17\,\unit{\second}$ & Automatic travel-time residual gating \\
        Local Magnitude Residual ($|\Delta M_{\mathrm{L}}|$) & $\le 0.25$ &
        $\le 0.25$ & Station-subset
        distance clipping ($r \le 30\,\unit{\kilo\meter}$) \\
        End-to-End Latency & $< 60\,\unit{\second}$ & $< 180\,\unit{\second}$ &
        Leonardo Booster streaming pipeline execution \\
        \bottomrule
      \end{tabularx}
      \begin{tablenotes}[flushleft]
        \footnotesize
      \item[a] Higher nominal FDR in Response Mode reflects genuine
        microseismic detection during intense aftershock sequences.
      \item Target \& Scope: Operational configuration matrix for
        deploying the ML pipeline in real-time regional observatories.
      \item What the reader should notice: The system balances rapid,
        high-precision public alerting in Surveillance Mode with
        exhaustive aftershock cataloging in Response Mode.
      \item Evidence Anchor: Benchmarking parameters validated in
        Chapter~\ref{chap:results}.
      \end{tablenotes}
    \end{threeparttable}
  \end{table}

  By combining PhaseNet (INSTANCE) picking, PyOcto 4D association, and
  NonLinLoc 1D relocation, regional seismological institutions can
  achieve automated catalog compilation with sub-kilometer epicentral
  accuracy, sub-second origin time fidelity, and $>96\%$ event
  completeness, setting a new benchmark for AI-driven operational
  seismology in Europe.
\end{landscape}
